\documentclass[12pt]{article}
\usepackage[utf8]{inputenc}
\usepackage[T1]{fontenc}
\usepackage{hyperref}
\title{Wangkumara Language Resources for K-12 Curriculum Development}
\author{Compiled by Claude Sonnet 5.5 (claude-sonnet-5-5) for Andr\'as Kornai}
\date{}
\begin{document}
\maketitle

\section{Introduction}
Wangkumara (ISO 639-3 xwk) is a Karnic, Pama-Nyungan language of south-western Queensland, associated with the Wilson and Bulloo River country. It is extinct or has no fluent speakers, so the goal is a heritage and revitalization curriculum. The search found very little material, and only a few reference sources are cited here. Sources disagree on the date of last speakers (ASJP gives 1990, Wikipedia about 2005). AIATSIS lists two entries for the language (L25 Bulloo River and L68 Wilson River), so documentation may be spread across both.

\section{Primary Educational Resources}
No primer, textbook or school programme for Wangkumara turned up in the searches.

\section{Linguistic and Reference Materials}
\begin{itemize}
\item McDonald and Wurm (1979), \emph{Basic Materials in Wangkumara (Galali): Grammar, Sentences and Vocabulary}, Pacific Linguistics B-65. The State Library of Queensland says it is available online. See \url{https://www.slq.qld.gov.au/blog/2019-international-year-indigenous-languages-word-week-week-six}.
\item Robertson, Carol (1984), \emph{Wangkumara grammar and dictionary}, Aboriginal Education Unit, Department of TAFE, Sydney. This is a bibliographic record only, and AIATSIS lists access as closed. See \url{https://glottolog.org/resource/reference/id/6977}.
\item AIATSIS Austlang entries: \url{https://collection.aiatsis.gov.au/austlang/language/L68} and \url{https://collection.aiatsis.gov.au/austlang/language/L25}.
\item ASJP wordlist, taken from McDonald and Wurm: \url{https://asjp.clld.org/languages/WANGKUMARA}.
\item Wikipedia overview: \url{https://en.wikipedia.org/wiki/Wilson_River_language}.
\end{itemize}

\section{Religious and Community Resources}
No scripture translations or community organisation sites specific to Wangkumara were found in the searches.

\section{Digital and Audio Resources}
No audio, video or software resources were found in the searches. Contacting AIATSIS and the State Library of Queensland about archival recordings would be a sensible next step.

\section{Curriculum Development Context}
The documented material is a single 1979 grammar and vocabulary, plus a 1984 grammar and dictionary of unconfirmed availability. A curriculum would be a small revitalization one, built from the wordlist and sentences and developed with descendant communities. Nothing in these searches shows existing school use.
\end{document}